% latexmk -pdflatex='xelatex %O %S' -pvc -pdf main.tex
\documentclass{../util/algebra}
\setenumerate[1]{label=(\alph*),topsep=0pt}
\setenumerate[2]{label=(\roman*),topsep=0pt}

\title{Vorlesungsskript zur Algebra \\ im Wintersemester 26/27}
\author{Felix Cherubini}

\makeindex

\begin{document}

\maketitle

\tableofcontents

\vspace{1cm}~\\

\subsection*{Zu diesem Skript}

Dieses Skript entsteht als Begleitmaterial zur Vorlesung ``Einführung in die Algebra'', die ich im Wintersemester 26/27 an der Universität Augsburg halte.
Zweck dieses Skripts ist es, den Zuhörern und mir selbst zur Erinnerung an die Vorlesungsinhalte zu dienen --
bei der Beschäftigung mit dem Thema ist es hilfreich in Lehrbücher zu schauen.
Ich orientiere mich teilweise an Siegfried Boschs Lehrburch ``Algebra'', das jedoch einem anderen Aufbau folgt.

Wer Fehler findet und diese korrigieren oder darauf aufmerksam machen will, kann das auf der github-Seite dieses Skripts machen:
\href{https://github.com/felixwellen/}{https://github.com/felixwellen/}.
Dort gibt es auch die Möglichkeit, sogenannte ``Issues'' anzulegen.
Hier könnten sie etwa darüber informieren, wenn sie eine Passage nicht verstehen oder einen Fehler gefunden haben.
Sie haben über sogenannte ``Pull requests'' die Möglichkeit, Fehler auch selbst zu korrigieren.

\subsection*{Ablauf des Übungsbetriebs}

Für Studierende im Bachelor Mathematik wird es im Anschluss an die Vorlesungszeit mündliche Prüfungen geben.
Für alle anderen Teilnehmer der Vorlesung handelt es sich um eine Portfolio-Prüfung. Diese besteht aus den folgenden Teilen:
\begin{enumerate}[(i)]
\item Klausur nach der Vorlesungszeit
\item Präsentation eigener Lösungen zu Übungsaufgaben in den Übungsgruppen
\end{enumerate}
Für den zweiten Punkt werden wir auf den wöchentlichen Übungsblättern Aufgaben auszeichnen, die für das Vortragen in den Übungsgruppen geeignet sind. Wer eine Lösung präsentieren möchte, kann sich in der jeweiligen Übung vor Beginn in eine Liste eintragen. Wer tatsächlich dann auch vortragen kann, wird vom Tutor ausgelost.

Die Abgabe von Lösungen zu den Übungsblättern ist zwar freiwillig, aber trotzdem dringend empfohlen. Die Abgabe erfolgt in den Übungsgruppen.

\section{Gruppen}

\subsection{Grundlegende Definitionen}

\begin{definition}
  Sei $M$ eine Menge. 
  \begin{enumerate}[(a)]
  \item Eine Abbildung $\_\cdot\_:M\times M \to M$ nennen wir auch \begriff{Verknüpfung}.
  \item Ein \begriff{Magma} ist eine Menge $M$ zusammen mit einer Verknüpfung. 
  \item Eine Verknüpfung $\_\cdot\_:M\times M\to M$ heißt \begriff{assoziativ}, wenn gilt:
    \begin{align*}
      \forall x,y,z\in M. (x\cdot y)\cdot z = x\cdot (y\cdot z)
    \end{align*}
  \item Ein Magma mit assoziativer Verknüpfung heißt \begriff{Halbgruppe}
  \item Sei $M$ ein Magma mit Verknüpfung $\_\cdot\_$. Ein Element $e\in M$ heißt \begriff{linksneutral}, wenn
    \begin{align*}
      \forall x\in M. e\cdot x=x
    \end{align*}
    und \begriff{rechtsneutral}, wenn 
    \begin{align*}
      \forall x\in M. x\cdot e=x
    \end{align*}
    gilt. $e$ heißt \begriff{neutrales Element}, wenn $e$ links- und rechtsneutral ist.
  \item Ein Magma $(M,\_\cdot\_)$ zusammen mit einem neutralen Element $e\in M$ heißt \begriff{Monoid}.
  \item Sei $M$ ein Monoid und $x,y\in M$. $y$ heißt \begriff{rechtsinvers} zu $x$, wenn $x\cdot y=e$ gilt. $y$ heißt \begriff{linksinvers} zu $x$, wenn $y\cdot x=e$ gilt. Wenn $y$ link- und rechtsinvers zu $x$ ist, sagen wir $y$ ist \begriff{invers} zu $x$.
  \item Ein Monoid $(M,\_\cdot\_,e)$ heißt \begriff{Gruppe}, wenn es zu jedem Element ein inverses gibt.
  \end{enumerate}
\end{definition}

\begin{bemerkung}
  \begin{enumerate}[(a)]
  \item 
  \end{enumerate}
\end{bemerkung}


\printindex

\printbibliography

\end{document}
